\chapter{Preliminaries for Single-robot Temporal Task Planning}
\label{chap:preliminaries}

% Target length: 10--12 pages.

%==============================
\section{System Abstraction and Task Specification}
\label{sec:system_abstraction_task_specification}

A single robot is considered in a bounded workspace together with a local
 temporal task. The purpose of this section is to introduce the formal objects
 used throughout the chapter, namely the robot motion abstraction and the task
 language interpreted over that abstraction. The development proceeds from the
 continuous robot dynamics to a finite transition model and then to the temporal
 specification imposed on its runs.

\subsection{Continuous robot model}
\label{sssec:continuous_robot_model}

Let $x(t) \in \mathbb{R}^{n}$ denote the robot state and
 $u(t) \in \mathbb{R}^{m}$ denote the control input. The robot dynamics are
 described by
\begin{equation}
\dot{x}(t) = f\bigl(x(t),u(t)\bigr),
\label{eq:continuous-dynamics}
\end{equation}
where $f : \mathbb{R}^{n} \times \mathbb{R}^{m} \to \mathbb{R}^{n}$ is
 assumed to be locally Lipschitz. Under this regularity condition, every
 admissible control signal induces a unique state trajectory for each initial
 condition. The development below does not depend on a particular robot model.
 The only structural requirement is the availability of feedback controllers
 that can realize transitions among suitable regions of the workspace.

\subsection{Finite transition abstraction}
\label{sssec:finite_transition_abstraction}

Let $\mathcal{W} \triangleq$ a bounded workspace in $\mathbb{R}^{d}$, where
 $d \geq 2$ denotes the dimension of the environment. A finite partition of
 $\mathcal{W}$ is introduced as
 $\mathcal{R} \triangleq \{R_{1},\ldots,R_{M}\}$, with
 $\bigcup_{i=1}^{M} R_i = \mathcal{W}$ and
 $R_i \cap R_j = \varnothing$ for all $i \neq j$. Each $R_i \subseteq
 \mathcal{W}$ is interpreted as a region of interest. To encode task-relevant
 facts over the workspace, let
 $\mathsf{AP} \triangleq
 \mathsf{AP}_{\mathrm{reg}} \cup \mathsf{AP}_{\mathrm{prop}}$, where
 $\mathsf{AP}_{\mathrm{reg}}$ contains region propositions and
 $\mathsf{AP}_{\mathrm{prop}}$ contains semantic propositions such as pickup
 locations, delivery locations, or surveillance targets. A labeling function
 $L : \mathcal{R} \to 2^{\mathsf{AP}}$ assigns to each region $R_i$ the set
 $L(R_i)$ of atomic propositions that hold whenever the robot occupies that
 region. In particular, $L(R_i)$ contains the region proposition associated with
 $R_i$. This abstraction is standard in formal methods and robot motion
 planning \cite{belta2007symbolic,kloetzer2008fully,baier2008principles}.

\begin{definition}[Weighted finite transition system]
The motion abstraction of the robot is the weighted finite transition system
\begin{equation}
\mathcal{T} \triangleq
\bigl(S,S_{0},\rightarrow,\mathsf{AP},L_{\mathcal{T}},w\bigr),
\label{eq:transition-system}
\end{equation}
where $S \triangleq \{s_{1},\ldots,s_{M}\}$ is the set of abstract states,
 $S_{0} \subseteq S$ is the set of admissible initial states,
 $\rightarrow \subseteq S \times S$ is the transition relation induced by the
 executable motions, $L_{\mathcal{T}} : S \to 2^{\mathsf{AP}}$ is the
 induced labeling map, and $w : S \times S \to \mathbb{R}_{\geq 0}$ assigns
 a nonnegative cost to each enabled transition. The correspondence between
 regions and states is one-to-one, and
 $L_{\mathcal{T}}(s_i) \triangleq L(R_i)$ for each state $s_i$ associated with
 region $R_i$.
\end{definition}

The transition relation in \eqref{eq:transition-system} is control-driven
 rather than purely geometric. In particular, $s_i \rightarrow s_j$ is enabled
 only when there exists an admissible feedback law that drives the robot from
 region $R_i$ to region $R_j$ in finite time while respecting the relevant
 motion constraints. An infinite run of the transition system in
 \eqref{eq:transition-system} is a sequence
 $\rho \triangleq s_0 s_1 s_2 \cdots$, where $s_0 \in S_0$ and
 $s_k \rightarrow s_{k+1}$ for all $k \in \mathbb{N}$. The run $\rho$
 induces the trace
\begin{equation}
\mathrm{trace}(\rho) \triangleq
L_{\mathcal{T}}(s_0)L_{\mathcal{T}}(s_1)L_{\mathcal{T}}(s_2)\cdots,
\label{eq:trace}
\end{equation}
where each symbol belongs to the alphabet $2^{\mathsf{AP}}$. The word in
 \eqref{eq:trace} is the object on which the temporal task is interpreted.

\subsection{Temporal task specification}
\label{sssec:temporal_task_specification}

The local task assigned to the robot is specified in linear temporal logic
 \cite{baier2008principles,fainekos2009temporal,kress2009temporal}. Let
 $\varphi$ denote an LTL formula generated by
\begin{equation}
\varphi ::= \top \mid p \mid \neg \varphi \mid
\varphi_1 \wedge \varphi_2 \mid \bigcirc \varphi \mid
\varphi_1 \,\mathsf{U}\, \varphi_2,
\label{eq:ltl-syntax}
\end{equation}
where $p \in \mathsf{AP}$, $\bigcirc$ denotes the next operator, and
 $\mathsf{U}$ denotes the until operator. The derived operators
 $\vee$, $\Diamond$, $\Box$, and $\Rightarrow$ retain their standard
 meanings. This syntax is expressive enough to capture reachability, repeated
 visitation, ordering constraints, and response rules.

\begin{definition}[Task satisfaction]
Let $\sigma \triangleq \sigma_0 \sigma_1 \sigma_2 \cdots$ be an infinite
 word over $2^{\mathsf{AP}}$. The satisfaction relation for the syntax in
 \eqref{eq:ltl-syntax} is defined recursively by
\begin{equation}
\begin{aligned}
(\sigma,k) \models p
&\iff p \in \sigma_k, \\
(\sigma,k) \models \neg \varphi
&\iff (\sigma,k) \not\models \varphi, \\
(\sigma,k) \models \varphi_1 \wedge \varphi_2
&\iff (\sigma,k) \models \varphi_1
\ \text{and} \ (\sigma,k) \models \varphi_2, \\
(\sigma,k) \models \bigcirc \varphi
&\iff (\sigma,k+1) \models \varphi, \\
(\sigma,k) \models \varphi_1 \,\mathsf{U}\, \varphi_2
&\iff \exists \ell \geq k \ \text{such that }
(\sigma,\ell) \models \varphi_2 \\
&\qquad \text{and } (\sigma,j) \models \varphi_1
\ \text{for all } k \leq j < \ell,
\end{aligned}
\label{eq:ltl-semantics}
\end{equation}
where $\sigma_k$ denotes the set of propositions true at position $k$. The
 semantics in \eqref{eq:ltl-semantics} are evaluated at the initial position by
 $\sigma \models \varphi$ if and only if $(\sigma,0) \models \varphi$.
 Accordingly, a run $\rho$ of the transition system in
 \eqref{eq:transition-system} satisfies $\varphi$ if and only if
 $\mathrm{trace}(\rho) \models \varphi$.
\end{definition}

The system abstraction in
 Subsubsection~\ref{sssec:finite_transition_abstraction} and the temporal
 semantics in Subsubsection~\ref{sssec:temporal_task_specification} lead to the
 symbolic planning problem
\begin{equation}
\text{find a run } \rho \text{ of } \mathcal{T}
\text{ such that } \rho \models \varphi.
\label{eq:planning-problem}
\end{equation}
The next section introduces the corresponding automata-theoretic constructions,
 explains how the required product model is computed, and then describes the
 algorithmic synthesis of a hybrid control strategy.

%==============================
\section{Automata-Based Synthesis and Hybrid Control}
\label{sec:automata_synthesis_hybrid_control}

Given the weighted transition system $\mathcal{T}$ in
 \eqref{eq:transition-system} and a local LTL formula $\varphi$, the planning
 problem in \eqref{eq:planning-problem} is converted into an
 automata-theoretic search problem. The presentation below follows the logic of
 this chapter: formal definitions are introduced first, then the relevant
 product model is computed, and finally the synthesis and execution algorithms
 are described. This construction follows the classical automata-theoretic
 approach to model checking and temporal-logic planning; see, for instance,
 \cite{baier2008principles,vardi1986automata,gastin2001fast,
 kloetzer2008fully,kress2009temporal}.

\subsection{Nondeterministic B\"uchi automaton}
\label{sssec:nba_definition}

\begin{definition}[Nondeterministic B\"uchi automaton]
The automaton associated with the LTL formula $\varphi$ is the nondeterministic
 B\"uchi automaton
\begin{equation}
\mathcal{B}_{\varphi} \triangleq
\bigl(Q,Q_{0},2^{\mathsf{AP}},\delta,F\bigr),
\label{eq:buchi-automaton}
\end{equation}
where $Q$ is the finite set of automaton states, $Q_{0} \subseteq Q$ is the
 set of initial states, $\delta : Q \times 2^{\mathsf{AP}} \to 2^{Q}$ is the
 transition map, and $F \subseteq Q$ is the accepting set.

An infinite run of $\mathcal{B}_{\varphi}$ over a word
 $\sigma = \sigma_0\sigma_1\sigma_2\cdots$ is a sequence
 $q_{0}q_{1}q_{2}\cdots$ such that $q_{0} \in Q_{0}$ and
 $q_{k+1} \in \delta(q_{k},\sigma_{k})$ for all $k \in \mathbb{N}$. Such a
 run is accepting when it visits $F$ infinitely often. The accepted language of
 the automaton in \eqref{eq:buchi-automaton} is denoted by
 $\mathcal{L}_{\omega}(\mathcal{B}_{\varphi})$ and coincides with the set of
 words satisfying $\varphi$ \cite{baier2008principles,vardi1986automata,
 gastin2001fast}. Therefore, the search for a run of $\mathcal{T}$ that
 satisfies $\varphi$ can be reformulated as the search for an accepting run of
 a suitable product automaton.
\end{definition}

\subsection{Computation of the product automaton}
\label{sssec:product_automaton_computation}

\begin{definition}[Weighted product B\"uchi automaton]
The weighted product B\"uchi automaton induced by
 $\mathcal{T}$ and $\mathcal{B}_{\varphi}$ is
\begin{equation}
\mathcal{P} \triangleq \mathcal{T} \otimes \mathcal{B}_{\varphi}
=
\bigl(X,X_{0},\rightarrow_{\mathcal{P}},F_{\mathcal{P}},w_{\mathcal{P}}\bigr),
\label{eq:product-automaton}
\end{equation}
where $X \triangleq S \times Q$ is the product state space,
 $X_{0} \triangleq S_{0} \times Q_{0}$ is the set of product initial states,
 and $F_{\mathcal{P}} \triangleq S \times F$ is the product accepting set.
 The transition relation of the product automaton in
 \eqref{eq:product-automaton} is specified by
\begin{equation}
(s,q) \rightarrow_{\mathcal{P}} (s',q')
\iff
\bigl(s \rightarrow s'\bigr)
\ \text{and} \
\bigl(q' \in \delta(q,L_{\mathcal{T}}(s))\bigr),
\label{eq:product-transition}
\end{equation}
where the first condition enforces motion feasibility in
 \eqref{eq:transition-system}, and the second condition updates the automaton
 state according to the current robot label. The transition cost is inherited
 from the transition system by
 $w_{\mathcal{P}}\bigl((s,q),(s',q')\bigr) \triangleq w(s,s')$ whenever the
 product transition in \eqref{eq:product-transition} is enabled.
\end{definition}

The product automaton in \eqref{eq:product-automaton} is computed by combining
 each abstract motion state with each automaton state, and then retaining only
 those pairs whose outgoing transitions satisfy \eqref{eq:product-transition}.
 In this way, the product graph records both the feasible motion evolution in
 $\mathcal{T}$ and the logical progress of the specification automaton.
 Consequently, accepting runs of $\mathcal{P}$ encode exactly the robot
 behaviors that satisfy the local task formula \cite{kloetzer2008fully,
 kress2009temporal,baier2008principles}.

\subsection{Plan-synthesis algorithm}
\label{sssec:plan_synthesis_algorithm}

Since an infinite accepting run is not convenient to manipulate directly, a
 prefix--suffix representation is adopted.

\begin{definition}[Prefix--suffix accepting run]
An accepting run of the product automaton is sought in the form
\begin{equation}
\varrho \triangleq
\varrho^{\mathrm{pre}}
\bigl(\varrho^{\mathrm{suf}}\bigr)^{\omega},
\label{eq:prefix-suffix-run}
\end{equation}
where $\varrho^{\mathrm{pre}}$ is a finite path from an initial product state
 to an accepting product state, and $\varrho^{\mathrm{suf}}$ is a finite cycle
 that starts from that accepting state and returns to it. More explicitly,
 $\varrho^{\mathrm{pre}} \triangleq x_{0}x_{1}\cdots x_{K}$ and
 $\varrho^{\mathrm{suf}} \triangleq x_{K}x_{K+1}\cdots x_{N}x_{K}$, where
 $x_{0} \in X_{0}$ and $x_{K} \in F_{\mathcal{P}}$. The representation in
 \eqref{eq:prefix-suffix-run} is sufficient for B\"uchi acceptance, because
 the suffix cycle revisits an accepting state infinitely often.
 \end{definition}

The quality of an accepting run is measured by
\begin{equation}
J(\varrho) \triangleq
\sum_{k=0}^{K-1} w_{\mathcal{P}}(x_{k},x_{k+1})
+
\gamma \sum_{k=K}^{N} w_{\mathcal{P}}(x_{k},x_{k+1}),
\label{eq:run-cost}
\end{equation}
where $\gamma \geq 0$ is a design parameter that balances the transient cost
 along the prefix and the steady-state cost along the suffix, with the
 convention $x_{N+1} \triangleq x_{K}$. The synthesis problem is therefore to
 find an accepting run of the form \eqref{eq:prefix-suffix-run} that minimizes
 \eqref{eq:run-cost}. This prefix--suffix optimization is closely related to
 the constructions used in optimal temporal-logic path planning; see
 \cite{smith2011optimal,schuppan2005shortest}.


Algorithmically, the synthesis procedure has three stages. First,
 $\mathcal{B}_{\varphi}$ is computed from the specification and then combined
 with $\mathcal{T}$ to obtain the product model \eqref{eq:product-automaton}.
 Second, for each candidate accepting state, one shortest path from an initial
 product state and one shortest accepting cycle are computed in the product
 graph. Third, the pair that minimizes the cost in \eqref{eq:run-cost} is
 selected, yielding the optimal accepting run
 $\varrho^{\star} \triangleq
 \varrho^{\star,\mathrm{pre}}\bigl(\varrho^{\star,\mathrm{suf}}\bigr)^{\omega}$.
 The corresponding discrete robot plan is obtained by projecting
 $\varrho^{\star}$ onto the transition-system component, namely
\begin{equation}
\rho^{\star} \triangleq
\mathrm{proj}_{S}\bigl(\varrho^{\star}\bigr)
=
s_{0}s_{1}s_{2}\cdots,
\label{eq:optimal-discrete-plan}
\end{equation}
where each consecutive pair $(s_{k},s_{k+1})$ is an enabled transition in
 $\mathcal{T}$. By \eqref{eq:product-transition}, the plan in
 \eqref{eq:optimal-discrete-plan} is feasible in the sense of
 \eqref{eq:transition-system}, and its trace satisfies $\varphi$.
 Algorithm~\ref{alg:optimal_accepting_run} summarizes the synthesis procedure
 described above.


 %==============================
 \begin{algorithm}[t]
\caption{Optimal accepting-run synthesis}
\label{alg:optimal_accepting_run}
\Input{Weighted transition system $\mathcal{T}$, LTL formula $\varphi$}
\Output{Optimal accepting run $\varrho^\star$ and projected plan $\rho^\star$}

Construct the B\"uchi automaton $\mathcal{B}_\varphi$ from $\varphi$\;
Construct the product automaton $\mathcal{P}=\mathcal{T}\otimes\mathcal{B}_\varphi$\;
Initialize the best cost as $+\infty$\;

\ForEach{accepting state $x_f \in F_{\mathcal{P}}$}{
  Compute a shortest path from some $x_0 \in X_0$ to $x_f$\;
  Compute a shortest cycle from $x_f$ back to itself\;
  Form the prefix--suffix run $\varrho$\;
  Evaluate its cost $J(\varrho)$ using \eqref{eq:run-cost}\;
  \If{$J(\varrho)$ is smaller than the current best cost}{
    Update the best run as $\varrho^\star \leftarrow \varrho$\;
  }
}
Project $\varrho^\star$ onto the transition-system component to obtain
$\rho^\star = \mathrm{proj}_S(\varrho^\star)$\;
\Return{$\varrho^\star,\rho^\star$}\;
 \end{algorithm}
 %==============================

\subsection{Hybrid execution algorithm}
\label{sssec:hybrid_execution_algorithm}

The final step is to realize the symbolic plan
 \eqref{eq:optimal-discrete-plan} on the continuous system
 \eqref{eq:continuous-dynamics}. Let
 $\kappa_{s_{k},s_{k+1}} : \mathbb{R}^{n} \to \mathbb{R}^{m}$ denote the
 feedback controller associated with the abstract transition from $s_{k}$ to
 $s_{k+1}$. The hybrid control law is then given by
\begin{equation}
u(t) = \kappa_{s_{k},s_{k+1}}\bigl(x(t)\bigr),
\label{eq:hybrid-control-law}
\end{equation}
where the controller corresponding to the current symbolic transition is
 activated until the region associated with $s_{k+1}$ is reached. At that
 event, the control law in \eqref{eq:hybrid-control-law} is switched to the
 controller associated with the next transition in
 \eqref{eq:optimal-discrete-plan}. In this way, the prefix of the optimal run
 is executed once, while its suffix is repeated indefinitely. The resulting
 closed-loop behavior follows the discrete accepting plan at the abstraction
 layer. Its formal correctness properties are established in
 Section~\ref{sec:correctness_complexity}.

%==============================
\section{Correctness and Complexity}
\label{sec:correctness_complexity}

This section summarizes the correctness guarantees of the synthesis procedure in
 Section~\ref{sec:automata_synthesis_hybrid_control} and analyzes its
 computational complexity. The first part concerns the logical correctness of
 the product construction and the hybrid execution law. The second part
 distinguishes between the complexity of on-the-fly single-robot synthesis and
 the scalability limitations that arise under direct multi-robot centralization.

\subsection{Correctness of symbolic synthesis}
\label{sssec:correctness_symbolic_synthesis}

The first correctness statement concerns the relation between the product
 automaton in \eqref{eq:product-automaton} and the original transition system
 in \eqref{eq:transition-system}.

\begin{lemma}
\label{lem:projection-correctness}
Let $\varrho$ be an accepting run of the product automaton $\mathcal{P}$ in
\eqref{eq:product-automaton}. Then its projection
$\rho \triangleq \mathrm{proj}_{S}(\varrho)$ is a run of $\mathcal{T}$ and
 satisfies $\rho \models \varphi$.
\end{lemma}

\begin{proof}
Let $\varrho \triangleq (s_{0},q_{0})(s_{1},q_{1})(s_{2},q_{2})\cdots$, where
$\varrho$ is accepting in $\mathcal{P}$. By the product transition rule in
\eqref{eq:product-transition}, every consecutive pair of product states
 satisfies
$(s_{k},q_{k}) \rightarrow_{\mathcal{P}} (s_{k+1},q_{k+1})$
only if $s_{k} \rightarrow s_{k+1}$. Therefore, the projected sequence
$\rho = s_{0}s_{1}s_{2}\cdots$ is a run of $\mathcal{T}$.

Moreover, the second condition in \eqref{eq:product-transition} implies that
 the automaton component evolves according to the trace generated by $\rho$,
 namely $q_{k+1} \in \delta\bigl(q_{k},L_{\mathcal{T}}(s_{k})\bigr)$.
 Since $\varrho$ is accepting, the automaton component visits
$F_{\mathcal{P}} = S \times F$ infinitely often. Consequently,
$q_{0}q_{1}q_{2}\cdots$ is an accepting run of $\mathcal{B}_{\varphi}$ over
 the word $\mathrm{trace}(\rho)$. Since the accepted language of
$\mathcal{B}_{\varphi}$ coincides with the set of words satisfying
$\varphi$, it follows that $\mathrm{trace}(\rho) \models \varphi$. The
 claim then follows from the trace-based satisfaction semantics introduced in
 Subsubsection~\ref{sssec:temporal_task_specification}.
\end{proof}

Lemma~\ref{lem:projection-correctness} shows that an accepting product run is
 sufficient to obtain a correct symbolic plan. The next result connects this
 symbolic plan to the continuous execution generated by the hybrid control law
 in Subsubsection~\ref{sssec:hybrid_execution_algorithm}.

\subsection{Correctness of hybrid execution}
\label{sssec:correctness_hybrid_execution}

\begin{lemma}
\label{lem:hybrid-execution-correctness}
Let $\rho^{\star}$ be the discrete plan in
\eqref{eq:optimal-discrete-plan} obtained from the optimal accepting run. Suppose
 that, for every enabled transition $s_{k} \rightarrow s_{k+1}$ appearing in
$\rho^{\star}$, the controller $\kappa_{s_{k},s_{k+1}}$ in
\eqref{eq:hybrid-control-law} drives the robot from the region associated with
$s_k$ to the region associated with $s_{k+1}$ in finite time and produces no
 spurious intermediate region switch. Then the closed-loop execution generates a
 symbolic word $\sigma_{x}$ satisfying $\sigma_{x} =
\mathrm{trace}(\rho^{\star})$, and consequently $\sigma_{x} \models
\varphi$.
\end{lemma}

\begin{proof}
By hypothesis, each controller $\kappa_{s_{k},s_{k+1}}$ realizes exactly the
 abstract transition from $s_k$ to $s_{k+1}$. Since the switching logic in
\eqref{eq:hybrid-control-law} activates the next controller only after the
 target region has been reached, the sequence of visited abstract states along
 the continuous execution coincides with the symbolic plan $\rho^{\star}$.
 Therefore, the word generated by the execution satisfies
$\sigma_{x} = \mathrm{trace}(\rho^{\star})$. By
Lemma~\ref{lem:projection-correctness}, the run $\rho^{\star}$ satisfies
$\rho^{\star} \models \varphi$. Applying the trace-based semantics then
 yields $\sigma_{x} \models \varphi$.
\end{proof}

Lemmas~\ref{lem:projection-correctness} and
\ref{lem:hybrid-execution-correctness} establish the logical correctness of the
 single-robot pipeline. In particular, if the local motion primitives correctly
 realize the abstract transitions, then the accepting run synthesized in the
 product automaton induces a continuous execution that satisfies the local task
 formula. This type of automata-based correctness argument is standard in formal
 methods and temporal-logic motion planning; see, for instance,
 \cite{baier2008principles,vardi1986automata,belta2007symbolic,
 fainekos2009temporal,kress2009temporal}.

\subsection{Complexity under on-the-fly construction}
\label{sssec:complexity_on_the_fly}

The complexity analysis is divided into two cases. The first concerns the
 single-robot synthesis problem under on-the-fly product construction. Let
$n_{\varphi} \triangleq |\varphi|$, where $n_{\varphi}$ denotes the size of
 the LTL formula. Standard translation from LTL to B\"uchi automata yields
 \cite{baier2008principles,gastin2001fast}.
\begin{equation}
|Q| = 2^{\mathcal{O}(n_{\varphi})}, \qquad
|\delta| = 2^{\mathcal{O}(n_{\varphi})},
\label{eq:buchi-size-bound}
\end{equation}
where $Q$ and $\delta$ are the state and transition sets of
$\mathcal{B}_{\varphi}$ in \eqref{eq:buchi-automaton}. Once
$\mathcal{T}$ and $\mathcal{B}_{\varphi}$ are fixed, the full product
 automaton in \eqref{eq:product-automaton} satisfies
$|X| = |S|\,|Q|$ and
$|\rightarrow_{\mathcal{P}}| =
\mathcal{O}\bigl(|\rightarrow|\,|\delta|\bigr)$, where $|S|$ and
$|\rightarrow|$ denote the number of states and transitions in $\mathcal{T}$.

Under on-the-fly construction, only the subset of product states and transitions
 requested by the graph search is generated. Let
$X_{\mathrm{vis}} \triangleq$ the set of visited product states during the
 search, where $X_{\mathrm{vis}} \subseteq X$, and let
$\rightarrow_{\mathrm{vis}} \triangleq
\{(x,x') \in \rightarrow_{\mathcal{P}} \mid x \in X_{\mathrm{vis}}\}$,
 where $\rightarrow_{\mathrm{vis}}$ is the induced set of generated product
 transitions. If shortest-path and shortest-cycle computations are performed
 from all candidate initial and accepting states, then the effective synthesis
 time is bounded by
\begin{equation}
T_{\mathrm{otf}} =
\mathcal{O}\Bigl(
\bigl(|X_{0}| + |F_{\mathrm{vis}}|\bigr)
|\rightarrow_{\mathrm{vis}}|
\log |X_{\mathrm{vis}}|
\Bigr),
\label{eq:otf-time}
\end{equation}
where $F_{\mathrm{vis}} \triangleq F_{\mathcal{P}} \cap X_{\mathrm{vis}}$
 denotes the accepting product states that are actually generated. The
 corresponding memory requirement is
$M_{\mathrm{otf}} =
\mathcal{O}\bigl(|X_{\mathrm{vis}}| + |\rightarrow_{\mathrm{vis}}|\bigr)$.

The worst-case asymptotic bound of the on-the-fly procedure coincides with that
 of full construction, since $X_{\mathrm{vis}} = X$ and
$\rightarrow_{\mathrm{vis}} = \rightarrow_{\mathcal{P}}$ when the search
 explores the entire product graph. However, in practice the on-the-fly
 strategy can be substantially cheaper, because many product states need not be
 generated if they are unreachable or irrelevant to the optimal accepting run.
 This distinction mirrors the discussion in the Chapter 3 draft on full versus
 on-the-fly construction.

\subsection{Multi-robot scalability}
\label{sssec:multi_robot_scalability}

The above complexity is still acceptable for a single robot when both the
 workspace abstraction and the specification size remain moderate. The situation
 changes significantly if the same method is applied directly to a system with
 multiple robots through a centralized team model. Consider $N$ robots with
 individual transition systems
$\mathcal{T}_{i} \triangleq
\bigl(S_{i},S_{i,0},\rightarrow_{i},\mathsf{AP}_{i},L_{i},w_{i}\bigr)$,
 where $i \in \{1,\ldots,N\}$. A direct joint abstraction leads to the team
 state space $S_{\mathrm{team}} \triangleq \prod_{i=1}^{N} S_{i}$, where each
 joint state records the simultaneous configuration of all robots. Consequently,
$|S_{\mathrm{team}}| = \prod_{i=1}^{N} |S_{i}|$, where the growth is
 multiplicative in the number of robots.

If a team-level temporal specification $\Phi$ is translated into a B\"uchi
 automaton $\mathcal{B}_{\Phi} \triangleq
\bigl(Q_{\Phi},Q_{\Phi,0},2^{\mathsf{AP}_{\mathrm{team}}},
\delta_{\Phi},F_{\Phi}\bigr)$, then the corresponding centralized product
 space becomes
\begin{equation}
|X_{\mathrm{team}}| =
|Q_{\Phi}| \prod_{i=1}^{N} |S_{i}|,
\label{eq:team-product-cardinality}
\end{equation}
where the blow-up arises from both the coupled motion state and the task
 automaton. The transition set also grows combinatorially, since every
 synchronized or asynchronously compatible joint move contributes a team-level
 transition. Therefore, both the graph-search cost and the memory demand inherit
 the multiplicative growth in \eqref{eq:team-product-cardinality}.

The implications are immediate. Even when the local models are modest, the
 direct centralized construction becomes computationally prohibitive as the
 number of robots increases. This observation explains why single-robot methods
 cannot be transferred naively to the multi-robot setting. Instead, the later
 chapters will exploit local specifications, compositional abstractions, and
 coordination mechanisms in order to avoid the full centralized blow-up in
 \eqref{eq:team-product-cardinality}. Related scalability issues have also been
 noted in the literature on multi-robot temporal-logic planning and deployment;
 see, for example, \cite{kloetzer2010automatic,chen2012formal,
 kloetzer2011multi,guo2015multi}.

%==============================
\section{Case Study}
\label{sec:case_study}

This section illustrates the proposed single-robot framework on a TIAGo-based
 simulation example. The implementation combines three components: a finite
 transition model of the robot motion, an LTL planner that synthesizes a
 satisfying symbolic plan, and a FlexBE execution layer that dispatches the
 corresponding navigation and manipulation behaviors. The example is adapted
 from the publicly available \texttt{tiago\_ltl\_flexbe} implementation,
 which builds on the P\_MAS\_TG planner and related LTL-based robotic
 planning tools \cite{package,LTLfirst,guo2013motion}.

\subsection{Workspace and abstract model}
\label{sssec:case_study_workspace_model}

The workspace consists of three regions of interest associated with navigation
 waypoints. Let $P_{\mathrm{cs}} \triangleq \{p_{1},p_{2},p_{3}\}$, where
$p_{1} = (0.57,-10.01,0.95)$,
$p_{2} = (-3.25,-0.31,0.82)$, and
$p_{3} = (-2.89,-9.01,0.52)$. The triplets denote planar position and heading.
 The associated region alphabet is
$\mathsf{AP}_{\mathrm{reg}}^{\mathrm{cs}} \triangleq
\{r_{1},r_{2},r_{3}\}$, where $r_i$ indicates that the robot occupies the
 region centered at $p_i$.

Besides navigation, the implementation supports discrete action labels through
 the TIAGo \texttt{play\_motion} interface. For the present example, only two
 actions are needed, namely
$\mathsf{AP}_{\mathrm{act}}^{\mathrm{cs}} \triangleq
\{\mathsf{pick},\mathsf{reach}\}$, where
$\mathsf{pick} \triangleq \texttt{pick\_from\_floor}$ and
$\mathsf{reach} \triangleq \texttt{reach\_max}$. The complete proposition
 alphabet is therefore
$\mathsf{AP}^{\mathrm{cs}} \triangleq
\mathsf{AP}_{\mathrm{reg}}^{\mathrm{cs}} \cup
\mathsf{AP}_{\mathrm{act}}^{\mathrm{cs}}$, which instantiates the general
 alphabet introduced in Section~\ref{sec:system_abstraction_task_specification}.

The motion abstraction is fully connected in the implementation. Hence, every
 region is reachable from every other region. The motion cost is proportional to
 the Euclidean waypoint distance and is given by
$w_{\mathrm{mot}}^{\mathrm{cs}}(s_i,s_j) \triangleq
0.1\bigl(\|(p_i)_x - (p_j)_x,\,(p_i)_y - (p_j)_y\|_2 + 10^{-3}\bigr)$,
 where $(p_i)_x$ and $(p_i)_y$ denote the planar coordinates of waypoint $p_i$.
 The action abstraction is represented by
$w_{\mathrm{act}}^{\mathrm{cs}}(\alpha) \triangleq 10$ for each
$\alpha \in \{\mathsf{pick},\mathsf{reach}\}$.

To account for both navigation and discrete actions, the single-robot planning
 model is taken as the motion-action transition system
\begin{equation}
\mathcal{T}_{\mathrm{cs}} \triangleq
\bigl(S_{\mathrm{cs}},S_{\mathrm{cs},0},
\rightarrow_{\mathrm{cs}},\mathsf{AP}^{\mathrm{cs}},
L_{\mathrm{cs}},w_{\mathrm{cs}}\bigr),
\label{eq:case-study-full-ts}
\end{equation}
where
$S_{\mathrm{cs}} \triangleq
S_{\mathrm{mot}}^{\mathrm{cs}} \times
\{\mathsf{none},\mathsf{pick},\mathsf{reach}\}$ and
$S_{\mathrm{cs},0} \triangleq \{(s_{1},\mathsf{none})\}$. Here,
$\mathsf{none}$ denotes the absence of a discrete action, and the robot starts
 from the waypoint associated with $r_1$. The labeling map satisfies
$L_{\mathrm{cs}}(s_i,\mathsf{none}) \triangleq \{r_i\}$ and
$L_{\mathrm{cs}}(s_i,\alpha) \triangleq \{r_i,\alpha\}$ for each
$\alpha \in \{\mathsf{pick},\mathsf{reach}\}$. Navigation transitions are
 of the form $(s_i,\beta) \rightarrow_{\mathrm{cs}} (s_j,\mathsf{none})$,
 where $s_i$ and $s_j$ are motion-adjacent states and
$\beta \in \{\mathsf{none},\mathsf{pick},\mathsf{reach}\}$, while action
 transitions are of the form
$(s_i,\mathsf{none}) \rightarrow_{\mathrm{cs}} (s_i,\alpha)$. Their weights
 are inherited from $w_{\mathrm{mot}}^{\mathrm{cs}}$ and
$w_{\mathrm{act}}^{\mathrm{cs}}$, respectively.

\subsection{Task specification and synthesized plan}
\label{sssec:case_study_task_plan}

The case-study task requires the robot to go to region $r_2$, execute the
 pickup action there, and then reach region $r_3$ and execute the reaching
 action. The corresponding LTL formula is
\begin{equation}
\varphi_{\mathrm{cs}} \triangleq
\Diamond \Bigl(
(r_{2} \wedge \mathsf{pick})
\wedge
\Diamond (r_{3} \wedge \mathsf{reach})
\Bigr),
\label{eq:case-study-task}
\end{equation}
where \eqref{eq:case-study-task} is the symbolic counterpart of the task
 expression used in the implementation. Applying the synthesis procedure from
 Section~\ref{sec:automata_synthesis_hybrid_control} to
$\mathcal{T}_{\mathrm{cs}}$ and \eqref{eq:case-study-task} yields a product
 automaton of the form \eqref{eq:product-automaton}, and the corresponding
 optimal accepting run is obtained in prefix--suffix form as in
 Subsubsection~\ref{sssec:plan_synthesis_algorithm}.

For the present instance, the optimal discrete prefix is
\begin{equation}
\rho_{\mathrm{cs}}^{\star,\mathrm{pre}} \triangleq
(s_{1},\mathsf{none})
(s_{2},\mathsf{none})
(s_{2},\mathsf{pick})
(s_{3},\mathsf{none})
(s_{3},\mathsf{reach})
(s_{3},\mathsf{none}),
\label{eq:case-study-prefix}
\end{equation}
where the first motion step drives the robot from $r_1$ to $r_2$, the second
 step executes the pickup action at $r_2$, the third step drives the robot from
$r_2$ to $r_3$, and the fourth step executes the reaching action at $r_3$. A
 trivial accepting suffix is given by
$\rho_{\mathrm{cs}}^{\star,\mathrm{suf}} \triangleq
(s_{3},\mathsf{none})^{\omega}$, where the robot remains in the terminal
 region after task completion. Since the motion graph is fully connected, the
 only nontrivial cost comparison concerns the ordering of the required visits
 and actions.

Using the motion and action costs defined in
Subsubsection~\ref{sssec:case_study_workspace_model}, the prefix cost induced
 by \eqref{eq:case-study-prefix} is
\begin{equation}
J_{\mathrm{cs}}^{\mathrm{pre}}
\triangleq
w_{\mathrm{mot}}^{\mathrm{cs}}(s_{1},s_{2})
+
w_{\mathrm{act}}^{\mathrm{cs}}(\mathsf{pick})
+
w_{\mathrm{mot}}^{\mathrm{cs}}(s_{2},s_{3})
+
w_{\mathrm{act}}^{\mathrm{cs}}(\mathsf{reach})
=
21.9134,
\label{eq:case-study-prefix-cost}
\end{equation}
where the numerical value follows from the waypoint coordinates. The plan in
 \eqref{eq:case-study-prefix} is therefore optimal for the instantiated model,
 since any feasible alternative would either preserve the same action cost while
 adding extra motion, or violate the ordering imposed by
 \eqref{eq:case-study-task}.

\subsection{Hybrid execution and discussion}
\label{sssec:case_study_execution_discussion}

The symbolic plan in \eqref{eq:case-study-prefix} is executed through the
 hybrid law introduced in \eqref{eq:hybrid-control-law}. In implementation
 terms, the continuous execution consists of the sequence
\begin{equation}
\texttt{goto}(r_{2})
\rightarrow
\texttt{pick\_from\_floor}
\rightarrow
\texttt{goto}(r_{3})
\rightarrow
\texttt{reach\_max},
\label{eq:case-study-execution-sequence}
\end{equation}
where \eqref{eq:case-study-execution-sequence} is the implementation-level
 counterpart of the discrete plan in \eqref{eq:case-study-prefix}. Under the
 assumptions of Lemma~\ref{lem:hybrid-execution-correctness}, the resulting
 execution satisfies \eqref{eq:case-study-task}.

This case study illustrates how the abstract constructions introduced in
 Sections~\ref{sec:system_abstraction_task_specification} and
 \ref{sec:automata_synthesis_hybrid_control} can be instantiated in a practical
 single-robot setting. The example also shows that, once the workspace and the
 available actions have been encoded as a finite transition system, complex
 navigation-and-action tasks can be expressed compactly in temporal logic and
 executed systematically through a hybrid planning-and-control architecture.

 %==============================
\section{Summary}
\label{sec:ch02_summary}
